\subsection{Separating the point germs by the nine tangent lines}

The following calculation gives additional geometric information about
Task C.3 directly in the abelian quotient. We work with the actual maps
on fundamental groups induced by regular invertible functions on $U$.
Their coefficient field is $K_0=\mathbb Q(\zeta)$; enlarging a field of
definition to contain $K_0$ is a restriction map of section sets, hence
is permitted in the germ colimit.

\begin{statement}{C.8}{Proved}{An exact tangent-line separation identity}
Let $\ell_d$ be a row vector defining the flex tangent line at $d\in D$.
The nine rows can be chosen as
\[
 (1,\zeta^a,0),\quad(0,1,\zeta^a),\quad(\zeta^a,0,1)
 \qquad (a=0,1,2).
\]
Set $P_d=\ell_d^*\ell_d/2$, a rank-one Hermitian projector on
$\mathbb C^3$. The real-linear map
\[
 \operatorname{Herm}_3(\mathbb C)\longrightarrow\mathbb R^9,
 \qquad H\longmapsto\bigl(\operatorname{tr}(P_d H)\bigr)_d
\]
is an isomorphism, and
\[
 \sum_d P_d=3I_3,\qquad
 \operatorname{tr}(P_dP_e)=
 \begin{cases}1&d=e,\\1/4&d\ne e.\end{cases}
 \tag{C.8.1}
\]
The kernel after deleting coordinate $d$ is the line
$\mathbb R(4P_d-I_3)$. A positive semidefinite matrix with the
eight undeleted measurements of $vv^*$, where $v\ne0$ and
$\ell_d(v)=0$, is exactly $vv^*$.
\end{statement}
\begin{proof}
Each row has norm squared two. The scalar product of two distinct
rows has absolute value one: for rows with the same two nonzero
positions this follows from $|1+\zeta|=|1+\zeta^2|=1$; for
different positions their supports meet in one coordinate.
Summing the projectors cancels the off-diagonal terms and gives
three on each diagonal entry. This proves (C.8.1). Their Gram
matrix is $(3I_9+J_9)/4$, with eigenvalues $3$ and $3/4$, so the
nine projectors are a real basis of the nine-dimensional
Hermitian space. Moreover
\[
 \operatorname{tr}\bigl(P_e(4P_d-I_3)\bigr)=3\delta_{de},
\]
which proves the assertion about the deleted coordinate.

Write a matrix with the prescribed eight measurements as
$vv^*+t(4P_d-I_3)$, $t\in\mathbb R$. The vector spanning
$\operatorname{im}P_d$ is orthogonal to $v$, and its quadratic
value is $3t$ times its norm squared. A nonzero vector orthogonal
to both this line and $v$ has quadratic value $-t$ times its
norm squared. Positive semidefiniteness therefore gives both
$t\ge0$ and $t\le0$. Thus $t=0$.
The exact program \texttt{code/c\_point\_separation.py} verifies
the rank, all 81 dual pairings, all nine tangent annihilations,
and the eigenvalues $(3,-1,-1)$ over $\mathbb Q(\sqrt{-3})$.
\end{proof}

\begin{statement}{C.13}{Proved}{The initial cuspidally central obstruction over every number field}
Let $L\subset\overline\Q$ be a number field and let
$\alpha\in H^1(G_L,V)$, with image $\beta\in H^1(G_L,T)$.
Orient $\bigwedge^2T\simeq\Z_3(1)$ by the commutator in
the once-punctured elliptic curve, and write $\omega$ for
this alternating pairing. For each $v\mid3$ define
\[
 q_v(\beta)=\tfrac12\operatorname{inv}_v
       \omega_*(\beta_v\smile\beta_v)\in\Z_3.
 \tag{C.13.1}
\]
Here $\operatorname{inv}_v:H^2_{\rm cts}(G_{L_v},\Z_3(1))
\xrightarrow{\sim}\Z_3$ is the inverse limit of the
normalized local Brauer invariant maps. Then
\[
 \alpha\in\operatorname{im}
   \bigl(\operatorname{Sec}_L(N^{\rm cc})\to H^1(G_L,V)\bigr)
 \quad\Longleftrightarrow\quad
             q_v(\beta)=0\text{ for every }v\mid3.
 \tag{C.13.2}
\]
The tuple satisfies $\sum_{v\mid3}q_v(\beta)=0$.
Every nonempty fiber is a torsor under
$H^1(G_L,\Z_3(1))=\varprojlim_n L^\times/(L^\times)^{3^n}$,
with no further stabilizer quotient. In particular, every
abelian section lifts to $N^{\rm cc}$ over a number field
having exactly one place above three, including $\Q$.
This statement makes no assertion that every tuple subject
to the displayed sum occurs as a self-pairing.
\par\smallskip\noindent\emph{Manuscript comparison: \S51.16,
folios 733v--736v, scans 172--178, formulas (371)--(379).}
\end{statement}
\begin{proof}
We first specify the genuine arithmetic action on the
once-punctured elliptic class-two quotient $C$. The
automorphism $[x:y:z]\mapsto[y:x:z]$ is elliptic inversion
for origin $O$. It preserves $E\setminus\{O\}$ and has
derivative $-1$ at $O$. Its comparison from the tangent
$-1$ to the tangent $1$ admits a $G_\Q$-fixed pro-three
path: the compatible $3^n$th roots of $-1$ can all be
chosen to be $-1$. Composition with this path gives a
genuine involution $j$ of the based pro-three group,
commuting with the arithmetic action. Applying the
comparison twice gives the constant path, as the product
of the two chosen roots is one at every level.

The integral class-two logarithm is available at three:
its group law is the Baker--Campbell--Hausdorff law with
the single denominator two. The central exact sequence
of Lie lattices is
\[
 0\longrightarrow\bigwedge^2T\longrightarrow\log C
       \longrightarrow T\longrightarrow0.
\]
The involution acts as $-1$ on $T$ and as $1$ on the
center. The operator $(1-j)/2$ therefore gives a
canonical lift of $T$ to its negative eigenspace.
An arithmetic automorphism commutes with $j$, so it
preserves this lift. In these coordinates
\[
 (t,z)(t',z')=(t+t',z+z'+\tfrac12\omega(t,t')),
 \qquad g(t,z)=(g t,\chi_3(g)z).
 \tag{C.13.3}
\]
This proves the claimed diagonalization by an actual
based morphism. In particular the central extension
obstruction for a $T$-cocycle $\beta$ is represented by
the continuous two-cocycle
\[
 (g,h)\longmapsto\tfrac12\omega(\beta(g),g\beta(h)).
 \tag{C.13.4}
\]

At a finite place $v\nmid3$, the curve has good reduction,
so $T$ is unramified. The two Frobenius eigenvalues have
complex absolute value $\sqrt{Q_v}$. Thus both
$T^{F_v=1}$ and $T(-1)^{F_v=1}$ vanish. The local
inertia sequence gives
\[
 H^1(G_{L_v},T)=T/(F_v-1)T,
\]
a finite group. This equality also follows at the
integral level by the finite-coefficient inverse limit;
the invariant towers are finite and have zero first
derived limit. The self-pairing of any element of a
finite group is torsion, whereas
$H^2_{\rm cts}(G_{L_v},\Z_3(1))=\Z_3$ is torsion-free.
Consequently (C.13.4) restricts to zero at every such
place. Its restrictions at archimedean places vanish
because three is odd.

For clarity, the required global injectivity retains
the topology. At each finite level the Kummer sequence
identifies $H^2(G_L,\mu_{3^n})$ with $\operatorname{Br}(L)[3^n]$.
The global Brauer invariant exact sequence, followed
by inverse limit, identifies its limit with
\[
 \left\{(b_v)_v\in\mathcal R_{v\ {
m finite}}(\Z_3):
                       \sum_v b_v=0\right\}.
 \tag{C.13.5}
\]
The sum converges because the family is restricted.
Conversely, reduction of such a tuple modulo $3^n$
has finite support and zero sum, and hence gives the
unique global Brauer class at that level; these
classes are compatible. Statement F.5 gives
$\varprojlim{}^1 H^1(G_L,\mu_{3^n})=0$, so Jannsen's
sequence identifies the continuous integral second
cohomology itself with (C.13.5). Equivalently one can
apply the same Kummer inverse-limit argument to the
completed multiplicative group. The local and global
Brauer invariant theorems are used in their stated
normalizations \cite{MilneCFT}.

The only possibly nonzero entries of (C.13.4) are
therefore precisely (C.13.1). Formula (C.13.5) proves
both the sum relation and the vanishing equivalence.
The arithmetic fiber product
$N^{\rm cc}=C\times_TV$ transports this equivalence
to (C.13.2). The central lifting fiber is acted on
transitively by $H^1(G_L,\Z_3(1))$. Its potential
stabilizer is the image of $H^0(G_L,V)$, which is
zero: a good Frobenius outside three has no eigenvalue
one on $V\otimes\Q_3$, since its boundary and elliptic
eigenvalues have absolute values $Q_v$ and
$\sqrt{Q_v}$. As $V$ is torsion-free, the integral
invariants vanish as well. This proves the assertion
about fibers. A tuple supported at one place and
having sum zero is zero, which proves the last claim.
\end{proof}

\begin{statement}{C.9}{Proved}{Injective point germs and their separation from cusp germs}
The canonical map
\[
 U(\overline{\mathbb Q})\longrightarrow
 \varinjlim_L\operatorname{Sec}_L(N/\gamma_3N)
\]
is injective and $G_{\mathbb Q}$-equivariant. Its image is
disjoint from all the cuspidal germs. Both assertions already
hold after projecting to the germ colimit of $H^1(G_L,V)$.
The smallest field of definition of a point germ is the field
of definition of the point.
\end{statement}
\begin{proof}
For distinct cusps $d,e$, the regular invertible function
$f_{d,e}=\ell_d/\ell_e$ on $U$ has divisor $3d-3e$ on $E$.
Its morphism $U\to\mathbb G_m$ gives a continuous map of
geometric pro-three groups $N\to\mathbb Z_3(1)$, factoring
through $V$. On the difference of the section classes of two
points $A,B\in U(L)$, the induced cohomology map is exactly
\[
 \kappa_L\bigl(f_{d,e}(A)/f_{d,e}(B)\bigr).
 \tag{C.9.1}
\]
This follows by pulling back the covers $z^{3^n}=t$ on
$\mathbb G_m$; their fibers are the compatible division roots
used to define the Kummer cocycle. Thus the comparison is a
map of actual section torsors, with the tangential origin
cancelling in the difference.

The kernel of $L^\times\to\varprojlim_n L^\times/L^{\times3^n}$
is its group of roots of unity of order prime to three.
Indeed every finite valuation of an infinitely three-divisible
element is zero. Dirichlet's unit theorem then reduces the
assertion to the finite roots-of-unity group. In particular
every complex absolute value of a kernel element is one.

Equality of two abelian germs is equality after restriction to
one finite number field containing $A,B$ and $K_0$. By (C.9.1)
the ratios of all these function values have absolute value one.
Fix a complex embedding and normalize vector representatives
$a,b$ of $A,B$ by $\ell_e(a)=\ell_e(b)=1$.
This is possible because the tangent at a flex meets $E$ only
at that cusp and $A,B\in U$. We obtain
$|\ell_d(a)|^2=|\ell_d(b)|^2$ for all nine $d$.
C.8 gives $aa^*=bb^*$, hence the same projective point $A=B$.

Now compare a point $A\in U$ with a section in the packet of
a cusp $d$. Choose $e\ne d$. For every $j\ne d$, the function
$f_{j,e}$ has order zero at $d$ and is a unit in its local ring.
The morphism on the tangential packet is consequently the
constant Kummer class of $f_{j,e}(d)$: extracting roots of a
local unit gives that constant fiber, and the varying tangent
parameter has exponent zero. Equality of abelian germs would
therefore give, after normalization at $\ell_e$,
\[
 |\ell_j(a)|^2=|\ell_j(d)|^2\quad(j\ne d).
\]
The cusp vector satisfies $\ell_d(d)=0$. The final assertion
of C.8 applied to the positive matrix $aa^*$ forces $A=d$,
contradicting $A\in U$. This proves disjointness for every
algebraic cusp and every completed Kummer parameter.

Equivariance follows from functoriality of the point section
and restriction. Injectivity identifies the stabilizer of a
point germ with the stabilizer of the point, which proves the
assertion about its smallest field of definition.
\end{proof}
